%\documentclass[article]{acmart}
%link to the TIIS manuscript submission: https://mc.manuscriptcentral.com/tiis 
\documentclass[manuscript,review,anonymous]{acmart}
%\documentclass[manuscript]{acmart}
\acmJournal{TIIS}
\setcounter{tocdepth}{3}
\usepackage{array}
\usepackage{booktabs}
\usepackage{pdfpages}
\usepackage[super]{nth}
\usepackage[nameinlink,noabbrev]{cleveref}
%%
%% \BibTeX command to typeset BibTeX logo in the docs
\AtBeginDocument{%
  \providecommand\BibTeX{{%
    Bib\TeX}}}
\setcopyright{acmlicensed}
\copyrightyear{2026}
\acmYear{2026}
\acmDOI{XXXXXXX.XXXXXXX}
%% These commands are for a PROCEEDINGS abstract or paper.
%% \acmJournal[TIIS]{}{FebruaryXXXXX}{XXXXX}
%%
%%  Uncomment \acmBooktitle if the title of the proceedings is different
%%  from ``Proceedings of ...''!
%%
%%\acmBooktitle{Woodstock '18: ACM Symposium on Neural Gaze Detection,
%%  June 03--05, 2018, Woodstock, NY}
%\acmISBN{978-1-4503-XXXX-X/2026/03}

%%
%% Submission ID.
%% Use this when submitting an article to a sponsored event. You'll
%% receive a unique submission ID from the organizers
%% of the event, and this ID should be used as the parameter to this command.
\usepackage[acronym]{glossaries}
\makeglossaries
\input{acronyms} % just loads definitions, doesn’t typeset them

%%
%% end of the preamble, start of the body of the document source.
\begin{document}
%%
%% The "title" command has an optional parameter,
%% allowing the author to define a "short title" to be used in page headers.

%\title{From Experts to Language Models: Enhancing Online Expert Elicitation with LLMs}

\title{From Experts to Language Models: Enhancing Online Expert Elicitation with LLMs}

\begin{abstract}


\end{abstract}

\begin{CCSXML}
<ccs2012>
 <concept>
  <concept_id>00000000.0000000.0000000</concept_id>
  <concept_desc>Do Not Use This Code, Generate the Correct Terms for Your Paper</concept_desc>
  <concept_significance>500</concept_significance>
 </concept>
 <concept>
  <concept_id>00000000.00000000.00000000</concept_id>
  <concept_desc>Do Not Use This Code, Generate the Correct Terms for Your Paper</concept_desc>
  <concept_significance>300</concept_significance>
 </concept>
 <concept>
  <concept_id>00000000.00000000.00000000</concept_id>
  <concept_desc>Do Not Use This Code, Generate the Correct Terms for Your Paper</concept_desc>
  <concept_significance>100</concept_significance>
 </concept>
 <concept>
  <concept_id>00000000.00000000.00000000</concept_id>
  <concept_desc>Do Not Use This Code, Generate the Correct Terms for Your Paper</concept_desc>
  <concept_significance>100</concept_significance>
 </concept>
</ccs2012>
\end{CCSXML}

%\ccsdesc[500]{Do Not Use This Code~Generate the Correct Terms for Your Paper}
%\ccsdesc[300]{Do Not Use This Code~Generate the Correct Terms for Your Paper}
%\ccsdesc{Do Not Use This Code~Generate the Correct Terms for Your Paper}
%\ccsdesc[100]{Do Not Use This Code~Generate the Correct Terms for Your Paper}

%%
%% Keywords. The author(s) should pick words that accurately describe
%% the work being presented. Separate the keywords with commas.
\keywords{Expert elicitation, Large language models, Judgment errors}
%% A "teaser" image appears between the author and affiliation
%% information and the body of the document, and typically spans the
%% page.
\received{05 February 2026}
\received[revised]{XX XX XXXX}
\received[accepted]{XX XX XXXX}

%% This command processes the author and affiliation and title
%% information and builds the first part of the formatted document.
\clearpage
\maketitle
\section{Introduction}



\section{State of the Art}


\section{Methods}
This section contains descriptions of the sampling procedure and details for the \gls{dc} experts, \gls{llm}, and \gls{dc} solutions. Also, the \gls{ee} questions, experimental setup, and simulation. Further, the section describes the chosen tests and analytical procedures for all experiments.

\subsection{Sampling}
Decisions on the number of experts for the \gls{ee} are sometimes determined by the desired expert mix, availability and willingness, as well as budget and time constraints \cite{machadoCCS2021, ohaganEE2006}. Among the studies that discuss the number of experts per study are Knol and Colleagues (2010), who propose an average of 6 per study \cite{knolEE2010}. Several publications recruited more than 6 participants, suggesting availability, and study objectives determine the number of experts \cite{vossosAdoption2022, glasgoEE2018, whistonEE2022, machadoCCS2021}. Hence, in our initial study \cite{anonymized2025_ShiftingCurrents}, we left it open and recruited experts based on availability. 

This section elaborates on the sample and sampling procedure of the \gls{dc} solutions, \gls{dc} experts, and the \gls{llm}. The studied \gls{dc} solutions are part of the solutions proposed in the \textit{Anonymized} project, which also provided ethical clearance for the \gls{ee} with \gls{dc} experts. The project focuses on introducing innovative approaches to power system implementation and on integrating \gls{dc} solutions into four key demonstration areas: ports, buildings, industry, and data centers \cite{Anonymized}. 

\subsubsection{Sampling the DC experts and LLMs}
The \gls{ee} with \gls{dc} involved 24 experts in two \gls{ee} sessions held between September and October, 2024. The first session involved 11 experts from within the \textit{Anonymized} project consortium. The other session involved 13 experts referred by colleagues in the \gls{dc} industry. The distribution of expertise of experts from both studies is presented in \Cref{tab:DC-experts-sample}. Email was used for initial communication and sharing of the consent forms with experts prior to performing the elicitation tasks. 

Sampled \gls{llm} as presented in \Cref{tab:llms-sample}, we used 4 general-purpose open-source \gls{llm} to contribute to the open-source community. To reduce computational requirements, the selected models have quantized model files below 10GB and are trained with 7B to 14B parameter counts as distributed via Ollama. The models are \gls{gemma}, \gls{llama}, \gls{mistral}, and \gls{phi4}. For each model, we generated 6 variants, 2 for each role type (i.e., generalist (2), normative (2), and subject-matter (2)), at temperatures ranging from 0.65 to 0.90, with a step size of 0.05. 

% \usepackage{booktabs}
\begin{table}[h!]
  \centering
  \caption{The DC Experts Sample}
  \label{tab:DC-experts-sample}
  \small
  \begin{tabular}{p{0.30\linewidth} r}
    \toprule
    Expertise & Count \\
    \midrule
    DC power electronics        & 11 \\
    Academia / society          & 4  \\
    Ports electrification       & 3  \\
    DC protection systems       & 2  \\
    Data-center electrification & 1  \\
    Buildings electrification   & 1  \\
    DC interoperability         & 1  \\
    AC/DC power systems         & 1  \\
    \bottomrule
  \end{tabular}
\end{table}


%%Side by Side table 
\begin{table}[h!]
\centering
\caption{The \gls{llm} Information}
\vspace{-1em}
\label{tab:llms-sample}
\small
\begin{tabular}{lcc}
\hline
Model & Quantized model file & Parameter count \\
\hline
\texttt{Phi-4} & 9.1 GB & 14.7B \\
Gemma3:12B & 8.1 GB & 12.2B \\
LlaMa-Pro & 4.7 GB & 8.36B \\
Mistral & 4.1 GB & 7.25B \\
\hline
\end{tabular}
\end{table}

\subsubsection{The DC solution}
A total of 11 out of 30 \gls{dc} solutions proposed in the \textit{Anonymized} project were used in \gls{ee} with \gls{dc} experts and \gls{llm}. The selection criterion for the solutions is their applicability across multiple demonstrators (i.e., demo cross-cutting aspect). The \gls{dc} solutions are: Smart and sustainable \gls{dc} cables, \gls{dc} connectors, static protection systems, semiconductor-based circuit breakers, protection \gls{dc} system design tool, \gls{dc}-\gls{dc} converters, LVAC-LVDC interlink converters, \gls{dc} measurement devices, \gls{dc} solution design tools, network design tool for \gls{dc} solutions, and solid-state circuit breakers. Detailed explanations on each are found in \cite{anonymized2025_ShiftingCurrents}.

\subsection{The EE questions}
The \gls{ee} questions used and presented in this study, both for \gls{dc} experts and \gls{llm}, are part of the \gls{ee} protocol detailed in our earlier publication \cite{anonymized2025_ShiftingCurrents}. The \gls{ee} questions were maintained between sessions with \gls{dc} experts and those with \gls{llm} to ensure comparability between them. The questions are feasibility, importance, and barrier selection presented in \Cref{tab:protocol-questions}. Each participant (both \gls{dc} experts and the \gls{llm}) answered each question based on the \gls{ee} setup. For feasibility, the focus was on \textit{overall feasibility}. As for the importance, each participant rated each of the 11 \gls{dc} solutions. As for the barrier selection, each presented a maximum of 5 barriers from a list of 11. 

\begin{table}[h!]
\centering
\caption{Elicitation Protocol Questions}
\vspace{-1em}
\label{tab:protocol-questions}
\small
\begin{tabular}{c p{7cm} p{5cm}}
\hline
SNo & Question & Options/Instructions \\
\hline
1. & How feasible is the use of DC in the previously described demonstrators? (overall ratings) & Not able to respond, not feasible, somewhat feasible, feasible, very feasible. \\
2. & How important is this solution? (For each DC solution) & Not able to respond, not important, somewhat important, important, very important. \\
3. & Barriers to adoption (a total of eleven barriers were presented for experts to select all that they perceive apply) & Select the top 5. \\
\hline
\end{tabular}
\end{table}

\subsection{The experimental setup and simulation}
The initial phase of the experimental setup involves data from two \gls{ee} sessions with \gls{dc} experts. Whereas the simulation experiments involve data only from \gls{llm}. Though the sessions with \gls{dc} experts were group-based \cite{anonymized2025_ShiftingCurrents}, the results presented in this article excluded the experts' discussion to ensure a smooth blend with \gls{llm}- generated simulation responses.
%
\Cref{fig:experimental_setup} depicts 4 experimental setups of this study: group-based \gls{ee}, zero-shot, context, and anchors.  

First, with \gls{dc} experts, the two sessions used the same methodology, tools, and procedure. The \glspl{ee} study was conducted in Microsoft Teams \cite{microsoftTeams}, with Slido for real-time question-and-answer sessions and initial summaries, which were viewed by all members and used by the elicitation facilitator to trigger discussion among experts \cite{slidoWebsite}. Read.AI was used for meeting transcription and recording \cite{readAI}, with the transcription focused more on discussion between experts and out of scope for this paper.

The second to fourth experiments involved the selected models described in \Cref{tab:llms-sample}. We prompted model variants using a Python-scripted, stateless Ollama REST API to ensure that prompts from earlier runs do not influence subsequent outputs. Ollama is an open-source platform that supports variant configurations via a model file \cite{Ollama}. To account for variability (i.e., stochastic nature) in generation, multiple runs are commonly used in studies that engage \gls{llm} \cite{azam2024multiple_runs, alampara2025multiple_runs}. We run 50 iterations for each question, both in zero-shot and in context. As for the anchor and list presentation, we run 10 iterations (i.e., per anchor arm: word, example, low-numeric, and high-numeric), leading to a total of 40 iterations per \gls{ee} question for comparability with zero-shot and context experiments. 
 
The second experiment is based on zero-shot learning (i.e., no-domain background information on the project, demonstration areas, and \gls{dc} solutions), as shown in \Cref{fig:prompt_zero-shot}. The third experiment on context aligns with the information provided to \gls{dc} experts, as obtained from the Read.AI transcription. Since \gls{llm} does not work well with \gls{asr} texts \cite{wang2024ASR, udagawa2024ASR}. The transcript was post-edited in ChatGPT-4.1, with punctuation restored, grammar corrected, and the text lightly paraphrased into clear text. The resulting paraphrased text was then checked and verified by the authors. The context included detailed background information on the project, demonstration areas, and \gls{dc} solutions, as shown in \Cref{fig:prompt_context}. 
Further experiments on anchored prompts (i.e., triggered prompts on statements concerning \gls{dc} microgrids) are presented in \Cref{fig:prompt_anchor}.    
 
\begin{figure}
    \centering
    \includegraphics[width=0.7\linewidth, trim=1cm 18cm 3cm 2.5cm, clip]{figures/experimental_setup.pdf}
    \caption{Experimental setup of DC experts and LLM simulations}
    \label{fig:experimental_setup}
\end{figure}

The analysis compares the experiments: 
\begin{itemize}
    \item The \gls{ee} with \gls{dc} experts, and simulated \gls{llm} at zero-shot, answering the first research question, \textit{RQ1}.
    \item Then, assessment of the effect of context on \gls{llm} responses (i.e., context vs. zero-shot) forms the second research question, \textit{RQ2}.
    \item Third, to identify the \gls{ee} judgmental errors by identifying the effect of anchors on \gls{llm} responses (i.e., anchors vs. context), forming the third research question, \textit{RQ3}.
    \item Furthermore, \textit{RQ1}, \textit{RQ2}, and \textit{RQ3} are all based on three \gls{ee} questions—feasibility, importance, and barriers.
\end{itemize} 

The \textit{feasibility question} is based on ratings on the overall feasibility of \gls{dc} solutions. Whereby, the important question was to rate each of the \gls{dc} solutions. The ratings for the feasibility and importance questions are on a 5-point scale, as indicated in \Cref{tab:protocol-questions}. Finally, for barrier selection, one must choose the 5 most perceived barriers to the adoption of \gls{dc} solutions from the 11 listed in \Cref{tab:barrier_dictionary}. In addition, the barriers in the figures have been drawn in accordance with the barrier codes listed in the barrier dictionary, \Cref{tab:barrier_dictionary}.

% Requires: \usepackage{booktabs}
\begin{table*}[h!]
  \caption{Barrier Dictionary with Corresponding Codes}
  \label{tab:barrier_dictionary}
  \centering
  \begin{tabular}{@{} l p{0.82\textwidth} @{}}
    \toprule
    \textbf{Code} & \textbf{Barrier Label} \\
    \midrule
    B1  & power losses, quality and safety issues \\
    B2  & reduced reliability in DC devices \\
    B3  & lack of use-cases in which DC is advantageous \\
    B4  & uncertain utility interaction (net metering, utility ownership, and agreed standards) \\
    B5  & lack of pilot projects \\
    B6  & public perception of DC and readiness to 'champion' installations from DC projects \\
    B7  & incompatibility of DC systems components \\
    B8  & misconception and lack of knowledge lead to lengthy/expensive design and permits \\
    B9  & lack of enough trained personnel in DC systems \\
    B10 & uncertain regulatory roadmap \\
    B11 & high costs of DC solutions \\
    \bottomrule
  \end{tabular}
\end{table*}

\subsection{Tests and analytical procedures}

\subsubsection{Zero-shot vs. DC experts}
The comparative assessment of feasibility and importance questions in \gls{llm} using zero-shot versus \gls{dc} experts' responses involves median values and the 50\% span (i.e., \gls{iqr}) of responses across \gls{llm} and the \gls{dc} experts. Regarding the barrier question, the selection rate was used to assess overlap between the top 5 barriers selected by \gls{llm} and those selected by \gls{dc} experts.

\subsubsection{Context vs. zero-shot, and anchors vs. context}
To understand the effect of one condition on the other, for ordinal ratings (i.e., feasibility and importance), we used a \gls{clmm}. The \gls{clmm} is a statistical test for analyzing ordered ordinal data and accounts for random effects between groups \cite{chen2020CLMM, peng2024CLMM}.  
In analyses, one condition (i.e., reference condition) is kept constant, whereby the effect of the second condition (i.e., contrast condition) is expressed relative to the reference condition reported as (\textit{CONTRAST} vs. \textit{REFERENCE}).

The formula for \gls{clmm} as in \Cref{eq:clmm}: 
\begin{equation}
\label{eq:clmm}
  \Pr(Y_i\le c)=\operatorname{logit}^{-1}\!\big(\tau_c-\eta_i-b_{p[i]}-u_{s[i]}\big),\quad c=0,1,2,3,\;\tau_0<\tau_1<\tau_2<\tau_3,  
\end{equation}

Where $Y_i$ denotes each variant's feasibility or importance ordinal rating values for each individual observation $i$. In 5-point ordinal categories ${0,1,2,3,4}$, maximum cuts are given by $K-1=4$, therefore $c$ is the cut (threshold) index is used in cumulative probabilities, $c=0,1,2,3$ and a maximum of the $\tau_c$ are the cut-points that define the category boundaries are $\tau_0<\tau_1<\tau_2<\tau_3$ estimated in increasing order across ordinal categories. 

The $\eta_i$ is the fixed-effects linear predictor that highlights the relevant coefficient for rating $i$. Also, the $b_p\sim\mathcal N(0,\sigma_p^2)$ and $u_s\sim\mathcal N(0,\sigma_s^2)$ are random intercepts per model and per \gls{dc} solution, respectively. Where, for feasibility, the ($u_s\equiv0$) as ratings are not per \gls{dc} solution. As for importance ratings, both random-intercept variables apply, as ratings differ across models and \gls{dc} solutions. 

The estimation of \gls{clmm} was performed using the 'R' ordinal package, and results were reported in two ways: odds ratios (ORs) and the probability of higher ratings. 
The OR for each model is derived from the logit link in \Cref{eq:clmm}. At [$OR = 1$] suggests no effect; [$ OR > 1$] the $CONTRAST$ is more likely to increase odds; [$ OR < 1$] more likely the $CONTRAST$ condition decreases odds. 
The probability of higher ratings of each model: $P_{\ge3}=\Pr(Y\ge3)$ is derived from fitted thresholds and $\eta_i$. It is computed per condition; therefore, a positive $\Delta$ indicates that the model is more likely to give higher ratings under the $CONTRAST$ condition than the $REFERENCE$ condition. Conversely, the chances are lower under the $CONTRAST$ condition.

\subsection{Methods for barrier selection}
As for the barrier list presentation question, the \gls{jsi} was computed. The \gls{jsi} is the measure of similarity between sets (irrespective of selection order), having a value ranging from $0$ to $1$ \cite{fender2017Jaccard, pereira2019Jaccard}. The closer its value to $1$, the more similar the two sets are, and the inverse is true.
We computed the \gls{jsi} for three conditions (i.e., context and zero-shot, context and anchor, and barrier lists) per model.   

For two finite sets $A$ and $B$, the \gls{jsi} is given by:

\begin{equation}
    J(A,B) \;=\; \frac{|A \cap B|}{|A \cup B|} \;\in\; [0,1].
\end{equation}

\section{Results}\label{sec:results}
This section presents the results of the four experiments presented as a comparative assessment, presenting: \gls{llm} at zero-shot versus \gls{dc} experts, the effect of context relative to zero-shot, and the effect of the anchors relative to context.

\subsection{DC experts vs. Models at zero-shot}\label{sec:zero-shot_DC-experts}
This part presents results of \gls{llm} at zero-shot with \gls{dc} experts on three questions: feasibility, importance, and barrier selection.

\subsubsection{Feasibility rating question—models vs. DC experts} \label{subsec:feas_zero_DCexperts}

The results on feasibility ratings are presented in \Cref{fig:rq1_feas_boxplots}, showing \gls{dc} experts and all models strongly assert on the feasibility of the \gls{dc} solutions, having [median=3]. All models have a tighter convergence [$IQR=0$] than a slightly looser one on \gls{dc} experts [$IQR=0.75$].
 
\begin{figure}
    \centering
    \includegraphics[width=0.75\linewidth]{figures/rq1_feas_zeroshot.png}
    \caption{Feasibility ratings—models vs. DC experts}
    \label{fig:rq1_feas_boxplots}
\end{figure} 

\subsubsection{Importance rating question—models vs. DC experts}\label{subsec:importance_zero_DCexperts}

The results of the importance ratings for each DC solution are summarized in box plots (\Cref{fig:rq1_importance_boxplots}) that compare responses from DC experts and the four LLMs. Overall, the models captured patterns broadly consistent with \gls{dc} experts, though differences emerged in the strength of their agreement on specific solutions.

For most DC solutions, the median importance rating across groups was 4 (“very important”), suggesting a shared understanding of the relative significance of these technologies. In nearly all cases, the models exhibited tighter convergence (IQR = 0) than the DC experts, who displayed slightly broader variation (IQR $\approx$ 0.5–1). This indicates that LLM responses were more homogeneous and less dispersed than those from human participants.

Some examples help illustrate these patterns:
\begin{itemize}
    \item For the DC–DC converter, both \gls{phi4} and \gls{gemma} closely mirrored the experts’ evaluation (median = 4, IQR = 0), showing nearly identical central tendencies and dispersion.
    \item The DC connectors question revealed similar alignment: \gls{gemma}, \gls{mistral}, and \gls{phi4} all matched the experts’ median (4) but were more consistent internally (IQR = 0 vs. experts’ IQR = 1).
    \item In contrast, the DC measurement device showed slightly wider divergence: \gls{llama} had perfect agreement within its own outputs (IQR = 0), while experts showed mild variability (IQR = 0.5), and the other models were more dispersed (IQR = 1).
    \item For the LVAC–LVDC interlink converter, experts and \gls{gemma} shared the same median rating (4), but the experts’ responses were more cohesive (IQR = 0) than the model’s (IQR = 1), suggesting human participants agreed more strongly on this solution’s importance.
\end{itemize}

Across all solutions, Phi-4 most frequently reproduced the experts’ median ratings (10 of 11 solutions), followed by \gls{gemma} (7 solutions), \gls{mistral} (4), and \gls{llama} (1). When considering both medians and IQRs together, \gls{phi4} and \gls{gemma} consistently emerged as the closest approximations to expert judgment, combining accurate central tendencies with stable consensus across multiple \gls{dc} solutions.

In summary, while no model perfectly replicated expert responses across all criteria, \gls{phi4} and \gls{gemma} demonstrated the strongest alignment, suggesting they capture the general structure of expert reasoning more reliably than the other models.

\begin{figure}
    \centering
    \includegraphics[width=0.85\linewidth]{figures/rq1_imp_zeroshot.png}
    \caption{Importance ratings—models vs. DC experts}
    \label{fig:rq1_importance_boxplots}
\end{figure}

\subsubsection{Barrier selection question—models vs. DC experts}\label{subsec:barriers_zero_DCexperts}

\Cref{fig:rq1_barrier_zero-shot} presents the barrier selection rate in two facets. First \gls{dc} experts with pooled-\gls{llm} (i.e., combined responses from all models), then per model. 
Considering the \gls{dc} experts with pooled-\gls{llm}: indicate \gls{llm} highly selected \nth{1} to \nth{5} barrier, contrary to \gls{dc} experts who highly selected B6, B8, B9, B10, and B11. 
Per model comparison relative to \gls{dc} experts' highly rated barriers, shows: \gls{gemma} also highly rated 3 — B8, B9, and B11; \gls{phi4} highly rated 1 — B10; none from \gls{mistral} or \gls{llama}.

\begin{figure}
\centering
   \includegraphics[width=0.85\linewidth]{figures/rq1_barrier_selection.png}
    \caption{Barrier selection—models vs. DC experts}
   \label{fig:rq1_barrier_zero-shot}
\end{figure} 

\subsection{Context vs. zero-shot} \label{sec:context_zero-shot}
This section presents results on the \gls{llm} effect under context condition relative to the zero-shot condition for three \gls{ee} questions (i.e., feasibility, importance, and barriers) for each model.

\subsubsection{Feasibility question—context vs zero-shot} \label{subsec:rq2_feas}

\Cref{tab:rq2_OR_feas} shows the effect of context compared to zero-shot for the feasibility question. Three models (i.e., \gls{phi4}, \gls{gemma}, and \gls{mistral}) had no effect, having [($OR$=1), at 95\%(0.852, 1.174)]. On the other hand, (i.e., \gls{llama}) shows insignificant increasing odds at context relative to zero-shot [($OR$=1.104), at 95\%CI(0.927, 1.314)].
Furthermore, \Cref{fig:rq2_feas_dumbbell} shows there are no changes in probabilities of higher ratings for \gls{gemma}, and \gls{mistral} between context and zero-shot conditions. As for \gls{phi4}, the likelihood of higher ratings is high under zero-shot relative to context [$\Delta=-0.078$], while for \gls{llama}, the probability increases in context relative to zero-shot [$\Delta= 0.132$].

\begin{table*}
  \caption{Feasibility ratings—context vs. zero-shot}
  \label{tab:rq2_OR_feas}
  \begin{tabular}{lccc}
    \toprule
    Model & OR & CI\_l95 & CI\_u95 \\
    \midrule
    % replace the body with:
    Gemma3:12B  & 1.000 & 0.852 & 1.174 \\
    LlaMa-Pro   & 1.104 & 0.927 & 1.314 \\
    Mistral     & 1.000 & 0.852 & 1.174 \\
    Phi4        & 1.000 & 0.852 & 1.174 \\
    \bottomrule
  \end{tabular}
\end{table*}

\begin{figure}
    \centering
    \includegraphics[width=0.8\linewidth]{figures/rq2_feas_dumbbell.png}
    \caption{Feasibility question - the $\Delta$ in probabilities of higher ratings.}    
    \caption*{\footnotesize The annotations are the $\Delta$ between each model's $P\ge3$ at context and zero-shot. For instance, \gls{mistral} had the same value at both conditions, hence $\Delta=0.000$}
    \label{fig:rq2_feas_dumbbell}
\end{figure}

\subsubsection{Importance question—context vs zero-shot} \label{subsec:rq2_importance}
This part also reports the effect in two aspects: the OR per model in \Cref{tab:rq2_OR_importance}, and the probability of higher ratings indicated by the $\Delta$ in probabilities per model per \gls{dc} solution in \Cref{fig:rq2_dumbbell_importance}.

The OR results, as in \Cref{tab:rq2_OR_importance}, indicate all models have a significant increase in OR under the context condition relative to the zero-shot condition: \gls{phi4} (OR=4.675, 95\%CI[3.680, 5.940]), \gls{llama} (OR=3.898, 95\%CI[3.194, 4.757]), \gls{mistral}  (OR=3.010, 95\%CI[2.474, 3.661]), and \gls{gemma}  (OR=2.574, 95\%CI[2.287, 2.895]). 

Overall, \Cref{fig:rq2_dumbbell_importance} shows that none of the models have a negative $\Delta$ in probabilities, which indicates a positive shift towards higher ratings in context compared to zero-shot for all models. 

\gls{llama} leads other models in having the highest $\Delta$ value in 8 \gls{dc} solutions. This implies that it has a higher chance of achieving higher ratings across the 8 DC solutions in context than with zero-shot. The \gls{dc} solutions that \gls{llama} lead are: DC-DC converters, DC connectors, DC solution design tool, network design tool for DC solutions, protection DC system design tool, semiconductor-based circuit breaker, smart and sustainable DC cables, and solid-state circuit breaker. \gls{mistral} leads in two DC solutions: the \gls{dc} measurement device and static protection system. \gls{phi4} leads to one DC solution: the LVAC-LVDC interlink converter, and \gls{gemma} to none. 

\begin{table*}
  \caption{Importance ratings—context vs. zero-shot}
  \label{tab:rq2_OR_importance}
  \begin{tabular}{lccc}
    \toprule
    Model & OR & CI\_l95 & CI\_u95 \\
    \midrule
    % replace the body with:
    Gemma3:12B  & 2.574 & 2.287 & 2.895 \\
    LlaMa-Pro   & 3.898 & 3.194 & 4.757 \\
    Mistral & 3.010 & 2.474 & 3.661 \\
    Phi4    & 4.675 & 3.680 & 5.940 \\
    \bottomrule
  \end{tabular}
\end{table*}

\begin{figure}
    \centering
    \includegraphics[width=\linewidth]{figures/rq2_importance_dumbbells.png}
    \caption{Importance question - the $\Delta$ in probabilities of higher ratings.}
    \caption*{\footnotesize The annotations in each sub-figure representing the DC solutions are only for a model with the highest $\Delta$ of probabilities between context and zero-shot conditions.}
    \label{fig:rq2_dumbbell_importance}
\end{figure}

\subsubsection{Barrier selection question—context vs zero-shot} \label{subsec:rq2_barriers}

\Cref{tab:rq2_jaccard_top5} presents the \gls{jsi} values and \Cref{fig:rq2_barrier_slopegraph} indicates a change in model preferences on the top-5 barriers between the conditions.

\gls{llama} and \gls{mistral} selected the same barriers under the context condition relative to zero-shot, [$JSI$=1.000]. The selection pattern is captured in \Cref{fig:rq2_barrier_slopegraph}, which shows that \gls{llama} maintained the same barrier positions, while \gls{mistral} flipped third and fourth positions between the same barriers (i.e., B2 and B5). Furthermore, there were observed changes in the top-5 barrier selection behavior on \gls{gemma} [$JSI$=0.667], and highest on \gls{phi4} [$JSI$=0.429]. 

\begin{table*}
  \caption{JSI on top-5 barriers—context vs. zero-shot}
  \label{tab:rq2_jaccard_top5}
  \begin{tabular}{l c }
    \toprule
    Model& $JSI(\text{Context},\text{Zero-shot})$ \\
    \midrule
    Gemma3:12B & 0.667 \\
    LLaMa-Pro  & 1.000 \\
    Mistral    & 1.000 \\
    Phi4       & 0.429 \\
    \bottomrule
  \end{tabular}
  \begin{minipage}{0.95\textwidth}\small
    The \gls{jsi} was computed for the 5 barriers with the highest selection rate in each model, across the context and zero-shot conditions.
  \end{minipage}
\end{table*}

\begin{figure}
    \centering
    \includegraphics[width=0.85\linewidth]{figures/rq2_barrier_slopegraph.png}
    \caption{Barrier selection—context vs zero-shot}
    \caption*{\footnotesize \textbf{} This figure captures the changes in barrier selection pattern for each model under the context relative to the zero-shot.}
    \label{fig:rq2_barrier_slopegraph}
\end{figure}

\subsection{Anchors vs. context}  \label{res:rq3}
The results presented in this section assess the availability of judgmental errors when prompts are anchored on four arms (i.e., word, example, low-numeric, and high-numeric) on three \gls{ee} questions as well. Three judgmental errors are examined in the results: overconfidence, over-extremity, and anchoring and adjustment.

\subsubsection{Feasibility rating question—anchors vs. context} \label{subsec:rq3_feas}

\Cref{tab:rq3_feas_or_anchors_vs_context} shows that two models significantly altered their responses per anchor arm relative to the context. All other models (i.e., \gls{phi4}, \gls{gemma}) had no effect in all arms relative to the context condition, [$OR = 1$], therefore were omitted from the table. 
\gls{llama} adjusted in each of the anchor arms; word, example, low numeric, and high numeric had significant increasing odds: [(OR=3.729), 95\% CI(2.005, 6.933)], [(OR=2.844), 95\% CI(1.479, 5.468)], [(OR=2.343), 95\% CI(1.202, 4.565)], and [(OR=4.118), 95\% CI(2.215, 7.654)] respectively. The observed rhythmic tone has a low effect on low-numeric and a higher effect on high-numeric, exhibiting the \textit{anchoring and adjustment} judgmental error.  
In addition, though at an insignificant effect, \gls{mistral} also adjusted to the word anchor [(OR=1.359), 95\% CI(0.412, 4.489)]. The effect, which also positively increases its $\Delta$ on probability of higher ratings [$\Delta=0.007$], even more than \gls{llama}'s [$\Delta=0.003$] at the word anchor as presented in \Cref{fig:rq3_feas_dumbbells}.

On the other hand, \gls{gemma} and \gls{phi4} exhibited \textit{overconfidence} and \textit{over-extremity} behaviors by not adjusting their ratings under any of the anchor arms relative to context, but also being too rigid, both having [$\Delta=0$] in each of the anchor arms. 

\begin{table*}
  \caption{Feasibility ratings—anchors vs. context}
  \label{tab:rq3_feas_or_anchors_vs_context}
  \centering
  \begin{tabular}{llccc}
    \toprule
    Model & Arm         & OR & CI\_l95 & CI\_u95 \\
    \midrule
    LlaMa-Pro & Word          & 3.729 & 2.005 & 6.933 \\
    LlaMa-Pro & Example       & 2.844 & 1.479 & 5.468 \\
    LlaMa-Pro & Low numeric   & 2.343 & 1.202 & 4.565 \\
    LlaMa-Pro & High numeric  & 4.118 & 2.215 & 7.654 \\
    Mistral   & Word          & 1.359 & 0.412 & 4.489 \\
    \bottomrule
  \end{tabular}
\end{table*}

\begin{figure}
    \centering
    \includegraphics[width=0.85\linewidth]{figures/rq3_feas_dumbbells.png}
    \caption{Feasibility question—$\Delta$ in probabilities of higher ratings}
    \label{fig:rq3_feas_dumbbells}
\end{figure}

\subsubsection{Importance rating question—anchors vs. context} \label{subsec:rq3_importance}
\Cref{tab:rq3_imp_OR} presents the effect per model for the importance question. 
As for \gls{llama}, the effect, though it is significant at high-numeric and not significant at low-numeric, anchoring and adjustments judgmental error is observed as in the feasibility question: [(OR=1.857), 95\% CI(1.468, 2.349)] at high-numeric, and [(OR=1.202), 95\% CI(0.952, 1.518)] at low-numeric. In addition, \gls{mistral} and \gls{phi4} behaved as \gls{llama}, though not with significant effect. 

On the other hand, for the word anchor \gls{phi4}, \gls{gemma}, and \gls{llama} had significantly higher increasing odds of: [(OR=10.681), 95\% CI(4.135, 27.590)], [(OR=1.721), 95\% CI(1.190, 2.489)], and [(OR=1.393), 95\% CI(1.103, 1.759)] respectively. 
Further analysis is on the probabilities of higher ratings for each model in each anchor relative to the context condition. Since in both conditions (i.e., context and each anchor arm) there are no ratings less than 3, as presented in \Cref{fig:rq3_imp_heatmap}, therefore, by the trivial, its value is $1$.

\begin{table*}
  \caption{Importance ratings—anchors vs. context}
  \label{tab:rq3_imp_OR}
  \centering
  \begin{tabular}{llccc}
    \toprule
    Model & Arm & OR & CI\_l95 & CI\_u95 \\
    \midrule
    Gemma3:12B & Word       & 1.721   & 1.190 & 2.489 \\
    LlaMa-Pro  & Word       & 1.393   & 1.103 & 1.759 \\
    Phi4       & Word       & 10.681  & 4.135 & 27.590 \\
    LlaMa-Pro  & High numeric & 1.857  & 1.468 & 2.349 \\
    Mistral	   & High numeric & 0.584  & 0.446	& 0.764 \\
    Phi4	   & High numeric & 1.322	& 0.815	& 2.143 \\
    LlaMa-Pro  & Low numeric  & 1.202	& 0.952	& 1.518 \\
    Mistral	   & Low numeric  & 0.118	& 0.089	& 0.157 \\
    Phi4	   & Low numeric  & 0.369	& 0.253	& 0.540 \\
    \bottomrule
  \end{tabular}
\end{table*}

\begin{figure}
    \centering
    \includegraphics[width=0.85\linewidth]{figures/rq3_imp_heatmap_scenario_x_rating.png}
    \caption{Importance rating distribution—anchors vs. context}
    \label{fig:rq3_imp_heatmap}
\end{figure}

\subsubsection{Barrier selection question—anchors vs context} \label{subsec:rq3_barriers}

\Cref{tab:rq3_jaccard_top5} indicates the \gls{jsi} in two aspects. First context with pooled anchors (i.e., treat ratings irrespective of anchor arm), then per anchor arm.  
Results indicate that at pooled anchors, there is no difference in how three models (i.e, \gls{phi4}, \gls{mistral}, and \gls{gemma}) selected the top 5 barriers, hence [$JSI=1$]. Only \gls{llama} altered its selection pattern, [$JSI=0.667$]. 

Per model, per anchor arm; \gls{gemma} has [$JSI=1$] at the context condition for each of the three arms (i.e, example, low-numeric, and high-numeric). All other models have a perfect index [$JSI=1$] in two different anchor arms on each: \gls{llama} on word and low-numeric, \gls{mistral} on low-numeric and high-numeric, and \gls{phi4} on word and high-numeric. All other remaining arms, for each model, have [$JSI=0.667$].

This implies that each model mostly maintained its selection of the top 5 barriers across anchor arms, relative to context. This implies the models showed little to no variation at anchors relative to context; therefore, they have exhibited overconfidence and over-extremity in their selection behavior.

\begin{table}[ht]
\centering
\caption{JSI on top-5 barriers—anchors vs. context}
\label{tab:rq3_jaccard_top5}
\begin{tabular}{lccccc}
\hline
Model & $JSI(\text{Ctx},\text{Anchors})$ & $JSI(\text{Ctx},\text{Word})$& $JSI(\text{Ctx},\text{Example})$& $JSI(\text{Ctx},\text{Low-num})$ & $JSI(\text{Ctx},\text{High-num})$\\ \hline
Gemma3:12B & 1.000 & 0.667 & 1.000 & 1.000 & 1.000 \\ 
LLaMa-Pro & 0.667 & 1.000 & 0.667 & 1.000 & 0.667 \\ 
Mistral & 1.000 & 0.667 & 0.667 & 1.000 & 1.000 \\
Phi4 & 1.000 & 1.000 & 0.667 & 0.667 & 1.000 \\ \hline
\end{tabular}
\end{table}

\section{Discussion}  \label{sec:discussion}
This section maps our findings from \Cref{sec:results} to the research questions and related literature. Finally, some implications for the design of EE systems that collaborate with Human-LLMs are also provided. 

\subsection{Answers to Research Questions}

\subsubsection{Models with \gls{dc} experts}
The experimental analyses of models at zero-shot with the \gls{dc} experts aimed to answer the first research question (i.e., \textit{RQ1}). The question aims to understand the capabilities of \gls{llm} to emulate experts. Based on the results obtained and presented at \Cref{sec:zero-shot_DC-experts}, it is reasonable to say that it is possible to emulate experts with \gls{llm}. However, several aspects should be considered, for instance, the varied behavior of model variants towards questions with similar rating scales, such as ordinal ratings for feasibility and importance questions. 

\subsubsection{Effect under context relative to zero-shot}
The experimental assessment on context with zero-shot targets to answer the second research question (i.e., \textit{RQ2}). It aims to understand the effect of context on \gls{llm} performance in \gls{ee} questions. The obtained results, presented in \Cref{sec:context_zero-shot}, show that there is an effect on \gls{llm} responses in the context condition relative to the zero-shot, although the effect is not consistent across all questions and models. While some models, like \gls{llama}, affected feasibility and importance questions, others, like \gls{gemma}, \gls{phi4}, and \gls{mistral}, only positively affected the importance question. On barrier selection patterns between the conditions, \gls{llama} and \gls{mistral} were rigid in selection patterns in the two conditions, while \gls{phi4} and \gls{gemma} adjusted. This implies that providing \gls{llm} with more specific context related to the subject under study can positively enhance their ratings. 

\subsubsection{Judgmental errors under anchors relative to context}
The analyses of the anchors and context experiments focus on answering the third research question (i.e., \textit{RQ3}). The question aims to identify the prevalence of judgment errors when \gls{ee} audiences are \gls{llm} rather than human experts. Our experiments with models reveal the likelihood of judgmental errors, which are reported to be possible when \gls{ee} involves human experts. As in other experiments reported in this study, no single model exhibited all three judgmental errors (i.e., overconfidence, over-extremity, and anchoring and adjustment). Our results highlight the need for caution in the use of anchors when dealing with \gls{llm} for \gls{ee} studies.

\subsection{Positioning Within Related Work}



\section{Conclusion, Limitations, and Future Works}
This section presents the study's conclusion, the observed limitations arising from the experimental setup, and the results. The section also concludes with a possible future extension to this work.

\subsection{Conclusion}



\subsection{Limitations}




\subsection{Future Works}

\section{Data availability}
The data and all necessary scripts and/or files to support the reproducibility of the experiments included in this manuscript have been made available in a GitHub repository: https://anonymous.4open.science/r/second-publication-6662/README.md. 
The repository contains the following scripts and/or files by categories:
\begin{itemize}
    \item The model files used to generate \gls{llm} variants varied by roles, and model type,
    \item Python scripts used for prompting the \gls{llm} variants, and,
    \item The data used for the analysis presented in this manuscript.
\end{itemize}   

\section{Declaration on the use of AI}
We used several AI services during the execution of various experiments and user studies for this paper. First \textit{Read.AI} for meeting recording and transcription with \gls{dc} experts. Then we used \textit{ChatGPT-4.1} for \gls{asr} text cleaning. The clean texts were used across two experiments, context-aware and anchored prompts.


\begin{acks}
We authors acknowledge the contribution of \gls{ee} participants both outside and within the project consortium. We also extend our appreciation to the funding from the \textit{Anonymized} project for the continued support during study execution.
\end{acks}

%%
%% The next two lines define the bibliography style to be used, and
%% the bibliography file.
\bibliographystyle{ACM-Reference-Format}
\bibliography{sample-base}

%%
%% If your work has an appendix, this is the place to put it.
\clearpage
\appendix

\section{Research Methods}
\subsection{Zero-shot}
\begin{figure}[h!]
    \centering
\includegraphics[width=0.8\linewidth, trim=1cm 21cm 1.2cm 0, clip]{figures/prompt_zero-shot.pdf}
    \caption{Zero-shot prompt structure}
    \caption*{\footnotesize \textbf{Note:} The figure demonstrates the structure of the prompt at the zero-shot condition. The \textit{preamble} section is in all questions. In addition, at \textit{dynamic block} it builds with \textit{list of DC solutions} if the question is either importance or feasibility; otherwise it builds with \textit{list of barriers}. The final prompt contains all the aforementioned blocks, along with the instructions and prompts for model variants.}
    \label{fig:prompt_zero-shot}
\end{figure}

\subsection{Context}
\begin{figure}[h!]
    \centering
    \includegraphics[width=0.8\linewidth, trim=1cm 16cm 1.2cm 0, clip]{figures/prompt_context.pdf}
    \caption{Context prompt structure}
    \caption*{\footnotesize \textbf{Note:} The figure demonstrates the structure of the prompt at the context condition. The \textit{preamble} section has two parts: the detailed context on the project and demonstration areas of the \cite{Anonymized} project. In addition, at \textit{dynamic block} it builds with \textit{detailed list of DC solutions} as in \Cref{app:dc-solutions} if the question is either importance or feasibility; otherwise it builds with \textit{list of barriers}. The final prompt contains all the aforementioned blocks, along with the instructions and prompts for model variants.}
    \label{fig:prompt_context}
\end{figure}

\subsection{Anchors}
\begin{figure}[h!]
    \centering
    \includegraphics[width=0.8\linewidth, trim=1cm 12cm 1.2cm 0, clip]
    {figures/prompt_anchor.pdf}
    %trim = left bottom right top
    \caption{Anchors prompt structure}
    \caption*{\footnotesize \textbf{Note:} The figure demonstrates the structure of the prompt at the anchor condition. It resembles \textit{context condition} at \textit{pre-amble} section and part of \textit{detailed list of DC solutions}. In addition, at \textit{dynamic block} the prompt alternates between the four anchor arms (i.e., word, example, low-numeric, and high-numeric). The use of either the DC solutions or \textit{list of barriers} depends on the types of \gls{ee} question. The final prompt contains all the aforementioned blocks for each question type, with the instructions. Then, variants of each model are prompted.}
    \label{fig:prompt_anchor}
\end{figure}

\end{document}
